\subsection{Synthetic sparse-coding data}
\label{sec:synthetic_sparse_coding_data}

We evaluated VG-SAE on a controlled sparse-coding model with known ground-truth
features, supports, and amplitudes. Let $d$ denote the observed activation
dimension and let $N$ denote the number of latent ground-truth features. A
dictionary $D^\star \in \mathbb{R}^{d \times N}$ was sampled column-wise and
normalized so that $\|d^\star_j\|_2 = 1$ for every feature $j$. To control
dictionary coherence, we first sampled independent Gaussian columns $u_j$ and a
shared Gaussian direction $c$, normalized $c$, and formed
\begin{equation}
    \tilde d_j
    =
    \sqrt{\kappa}\,c + \sqrt{1-\kappa}\,u_j ,
    \qquad
    d^\star_j = \frac{\tilde d_j}{\|\tilde d_j\|_2},
\end{equation}
where $\kappa \in [0,1)$ is a common-mode coherence parameter. When
$\kappa=0$, the dictionary columns are independent random directions up to
finite-dimensional correlations.

Feature activation probabilities were set by a density parameter $\rho$ and an
optional frequency-skew exponent $\alpha$. Specifically,
\begin{equation}
    q_j
    =
    \operatorname{clip}\left(
        \rho \frac{j^{-\alpha}}{N^{-1}\sum_{\ell=1}^{N}\ell^{-\alpha}},
        0,
        0.95
    \right),
    \qquad j=1,\ldots,N .
\end{equation}
For each sample $\mu$, the binary support variables were drawn independently as
$s^\star_{\mu j} \sim \mathrm{Bernoulli}(q_j)$. Conditional on the support,
nonnegative amplitudes were sampled as
\begin{equation}
    a^\star_{\mu j} \sim \mathrm{Exponential}(A),
    \qquad
    z^\star_{\mu j}=s^\star_{\mu j}a^\star_{\mu j},
\end{equation}
where $A$ is the amplitude scale. The clean observation was
\begin{equation}
    x^\star_\mu = D^\star z^\star_\mu ,
\end{equation}
and the observed input was generated by additive isotropic Gaussian noise,
\begin{equation}
    x_\mu = x^\star_\mu + \epsilon_\mu,
    \qquad
    \epsilon_\mu \sim \mathcal{N}(0,\sigma^2 I_d).
\end{equation}
The tensors $(x_\mu, z^\star_\mu, s^\star_\mu, D^\star)$ are retained, allowing
direct evaluation of reconstruction, dictionary recovery, amplitude shrinkage,
and support-selection quality.

\subsection{VG-SAE}
\label{sec:vg_sae_method}

The VG-SAE objective is an amortized, vector-output extension of the
Variational Garrote free energy. In the original VG sparse-regression model,
the output $y_\mu$ is scalar and each binary selector $s_i$ controls whether the
scalar regressor $w_i x_{\mu i}$ contributes to the prediction. In VG-SAE, each
sample has a vector observation $x_\mu \in \mathbb{R}^d$, and the learned
decoder column $d_j \in \mathbb{R}^d$ plays the role of a vector-valued
regression weight. The selector is sample-specific: $s_{\mu j}=1$ indicates
that latent feature $j$ is active for sample $\mu$.

The encoder parameterizes a mean-field Bernoulli posterior over support
variables and a deterministic amplitude:
\begin{align}
    \bar x_\mu &= x_\mu - b, \\
    m_{\mu j} &= q_\phi(s_{\mu j}=1\mid x_\mu)
              = \sigma(g_{\phi,j}(\bar x_\mu)), \\
    a_{\mu j} &= \operatorname{softplus}(r_{\phi,j}(\bar x_\mu)), \\
    h_{\mu j} &= m_{\mu j} a_{\mu j}.
\end{align}
Here $b$ is a trainable pre-bias, $g_\phi$ is the gate encoder, $r_\phi$ is the
amplitude encoder, and $h_\mu$ is the mean latent code. Decoder columns are
renormalized during training so that $\|d_j\|_2=1$. The mean reconstruction is
\begin{equation}
    \hat x_\mu = b + \sum_{j=1}^{N} m_{\mu j} a_{\mu j} d_j .
\end{equation}

Assuming isotropic Gaussian reconstruction noise, the scalar squared residual
in the VG paper is replaced by a squared Euclidean norm. Taking the expectation
over the mean-field Bernoulli posterior gives the per-sample energy
\begin{equation}
    E_\mu
    =
    \frac{1}{2}
    \left\|
        \bar x_\mu
        -
        \sum_{j=1}^{N}m_{\mu j}a_{\mu j}d_j
    \right\|_2^2
    +
    \frac{1}{2}
    \sum_{j=1}^{N}
        m_{\mu j}(1-m_{\mu j})a_{\mu j}^2\|d_j\|_2^2 .
    \label{eq:vg_sae_energy}
\end{equation}
The second term is the Bernoulli variance correction and follows from
$s_{\mu j}^2=s_{\mu j}$ and independence of the mean-field selectors.

The sparsity prior follows the VG paper:
\begin{equation}
    P(s_{\mu j}) =
    \frac{\exp(-\gamma s_{\mu j})}{1+\exp(-\gamma)},
    \qquad \gamma \in \mathbb{R} .
\end{equation}
Thus larger positive $\gamma$ penalizes active support variables, while
negative $\gamma$ defines a valid dense prior. The normalized fixed-$\beta$
VG-SAE objective over a
minibatch $\mathcal{B}$ is
\begin{equation}
    \mathcal{L}
    =
    \frac{1}{|\mathcal{B}|}
    \sum_{\mu \in \mathcal{B}}
    \left[
        \beta E_\mu
        - \frac{d}{2}\log\frac{\beta}{2\pi}
        +
        \gamma \sum_{j=1}^{N} m_{\mu j}
        + N\log(1+e^{-\gamma})
        -
        \eta \sum_{j=1}^{N} H(m_{\mu j})
    \right],
    \label{eq:vg_sae_fixed_beta}
\end{equation}
where
\begin{equation}
    H(m)=-m\log m-(1-m)\log(1-m)
\end{equation}
is the Bernoulli entropy and $\eta$ is an entropy coefficient. The default
experiments use the beta-eliminated form analogous to the VG paper. For a
minibatch containing $|\mathcal{B}|d$ scalar Gaussian observations,
\begin{equation}
    \beta^\star =
    \frac{|\mathcal{B}|d}{2\sum_{\mu \in \mathcal{B}} E_\mu},
\end{equation}
which yields the reduced objective
\begin{equation}
    \mathcal{L}_{\mathrm{red}}
    =
    \frac{1}{|\mathcal{B}|}
    \left[
        \frac{|\mathcal{B}|d}{2}
        \log\left(
            \frac{2\sum_{\mu \in \mathcal{B}}E_\mu}{|\mathcal{B}|d}
        \right)
        +
        \gamma
        \sum_{\mu \in \mathcal{B}}\sum_{j=1}^{N}m_{\mu j}
        + |\mathcal{B}|N\log(1+e^{-\gamma})
        -
        \eta
        \sum_{\mu \in \mathcal{B}}\sum_{j=1}^{N}H(m_{\mu j})
    \right].
    \label{eq:vg_sae_reduced}
\end{equation}
This is the direct vector-output analogue of the VG free energy: the scalar
sample count $M$ in the paper is replaced by the number of scalar reconstructed
coordinates $|\mathcal{B}|d$.

We implement three precision modes.  The profiled mode uses
Eq.~\eqref{eq:vg_sae_reduced}; the fixed mode uses
Eq.~\eqref{eq:vg_sae_fixed_beta}; and the learned mode parameterizes
$\beta=\exp(\log\beta)$ and optimizes the same full objective.  In the learned
mode,
\begin{equation}
  \frac{\partial \mathcal{L}}{\partial\log\beta}
  = \frac{\beta}{|\mathcal{B}|}\sum_{\mu\in\mathcal{B}}E_\mu-\frac{d}{2},
\end{equation}
so stationarity recovers the profiled minibatch value.  Entropy weights other
than one, or disabling entropy, are treated as ablations.

\subsection{SAELens-aligned baselines}
\label{sec:saelens_baselines}

All baselines are the official SAELens v6.47.0 (commit \texttt{8be1408})
training classes.  Their common path encodes $x-b_{\mathrm{dec}}$, adds
$b_{\mathrm{dec}}$ after decoding, and initializes the encoder from the decoder
transpose; decoder-norm handling remains architecture-specific.  Their
reconstruction loss is the minibatch mean squared Euclidean norm (SSE), without
the factor $1/2$ used in the VG energy.

The Standard/L1 SAE emits
$h=\operatorname{ReLU}(W_{\mathrm{enc}}(x-b_{\mathrm{dec}})+b_{\mathrm{enc}})$.
Its sparsity term is the configured coefficient times the minibatch mean of the
$\ell_p$ norm of $h_j\lVert d_j\rVert_2$ (with $p=1$ here).  Decoder rows are
not constrained to unit norm during Standard training.  At inference export,
the official norm-folding transform moves each nonzero decoder-row norm into
the corresponding latent scale, producing unit-norm decoder atoms while
preserving the reconstructed function.

TopK selects the $k$ largest raw encoder preactivations per sample and applies
ReLU only to the selected values.  Its AuxK loss asks up to $d/2$ dead features
to reconstruct the detached main-path residual, using the same decoder-norm
rescaling as the main TopK path and no decoder bias.  BatchTopK instead selects
$\lfloor k|\mathcal{B}|\rfloor$ positive activations globally during training.
It maintains a double-precision exponential moving average of the smallest
strictly positive selected activation and uses that fixed threshold at
evaluation, making inference independent of evaluation batch composition.

JumpReLU learns a positive per-feature threshold
$\theta_j=\exp(\log\theta_j)$ and emits
$h_j=z_j\mathbf{1}[z_j>\theta_j]$.  Threshold gradients use the rectangular
kernel estimator with bandwidth $0.05$; the sparsity loss is either the
corresponding step-function $L_0$ estimator or the optional tanh surrogate.
The Gated SAE shares encoder weights between
$\pi_{\mathrm{gate}}=W(x-b_{\mathrm{dec}})+b_{\mathrm{gate}}$ and the magnitude
path $W\exp(r_{\mathrm{mag}})(x-b_{\mathrm{dec}})+b_{\mathrm{mag}}$.  Its code
is the ReLU magnitude multiplied by the hard gate
$\mathbf{1}[\pi_{\mathrm{gate}}>0]$; training adds the decoder-norm-weighted
$L_1$ penalty on $\operatorname{ReLU}(\pi_{\mathrm{gate}})$ and a gate-only
auxiliary reconstruction SSE.

\subsection{Evaluation metrics}
\label{sec:vg_sae_metrics}

We report the mean reconstruction MSE, the model support density
\begin{equation}
    \rho_{\mathrm{model}}
    =
    \frac{1}{|\mathcal{B}|N}
    \sum_{\mu \in \mathcal{B}}\sum_{j=1}^{N}m_{\mu j},
\end{equation}
and the Bernoulli uncertainty
\begin{equation}
    V_{\mathrm{eff}}
    =
    \frac{1}{|\mathcal{B}|N}
    \sum_{\mu \in \mathcal{B}}\sum_{j=1}^{N}
    m_{\mu j}(1-m_{\mu j}).
\end{equation}
For synthetic data, learned decoder atoms are matched to ground-truth atoms by
maximum absolute cosine similarity using a linear assignment. After this
matching, support precision, support recall, decoder recovery cosine, and
amplitude shrinkage are computed against the known
$(s^\star_{\mu j}, z^\star_{\mu j}, D^\star)$.
